\section{Data and empirical design}\label{sec:short_design}

\subsection{WTI Average Price Options}

The empirical sample is built from individual Barchart histories for WTI Average Price Options. The raw archive contains 213 unique contracts and 11,179 contract-date observations, with expiry months ranging from September 2026 to June 2029. We compare model prices with the Barchart end-of-day settlement field. Calls and puts are retained as available, and moneyness is recomputed by valuation date rather than assigned once at the contract level.

For every observation, we identify the fixings already known and the futures contracts used for the remaining fixings. We initialize those contracts with the futures prices observed on the valuation date. This matters because the first-nearby contract can change during the averaging month.

Historical volatility is estimated from WTI futures returns available by the valuation date. The baseline uses the labelled continuous/front-month \texttt{CL=F} return series as the physical-measure proxy. A robustness check reconstructs a first-nearby return history contract by contract from CL settlements and excludes returns spanning contract switches.

\subsection{Out-of-sample volatility comparisons}

We evaluate volatility inputs through out-of-sample APO pricing errors. Implied volatility is estimated exclusively from option observations preceding each target date. Repricing uses the target-date futures curve, realized fixings, and discount factor; the volatility input alone comes from earlier option dates. Historical-volatility estimates follow the specified return-data cutoff. The sample comprises 2,381 contract-date observations across seven expiries and 15 target dates.

The baseline integrates conditional prices over a volatility posterior estimated from an expanding return sample. Historical alternatives use rolling windows of 63, 126, or 252 returns, an exponentially weighted moving average (EWMA) with a 63-business-day half-life, and a horizon-matched Gaussian GARCH(1,1) forecast. For GARCH, parameters are estimated from earlier returns; the target-date return updates conditional variance, which is then forecast over the remaining fixings. We map these forecasts to the constant annualized volatility that matches the remaining average's variance under the one-factor pricing model. Full comparisons are in Online Supplement Section S1 and Table S1.

We obtain contract-level APO implied volatilities by inverting the Curran pricing approximation. For each expiry, the previous-day implied-volatility smile is estimated on the latest available earlier date. A second specification fits a smile using all earlier dates, giving more weight to recent observations. Both depend on moneyness and option type. Neither uses the target-date APO settlement to estimate volatility.

To assess the value of smile fitting, we also carry forward each contract's most recent implied volatility. We reprice with the valuation-date curve, discount factor, and fixing schedule. This same-contract benchmark covers 2,366 observations.

Pricing performance is summarized by mean error, mean absolute error (MAE), and root mean squared error (RMSE). Dependence across contracts observed on the same valuation date is preserved with valuation-date cluster bootstrap resampling.

\subsection{Cross-instrument validation}

The within-APO comparison draws implied volatility from earlier observations of the same option family. We therefore add a cross-instrument check using standard monthly American WTI futures options (CME product LO) from Databento. We match LO option settlements to the underlying CL futures and invert prices using a Cox--Ross--Rubinstein futures-option tree.

For each APO target date, only LO observations from earlier dates are used. Implied-volatility smiles are fitted separately for the underlying CL futures contracts needed by the October-2026 APO fixing schedule. The resulting maturity-specific volatilities are mapped into the maintained scalar-volatility APO pricing model. We compare 253 APO contract-date observations over 10 target dates for which both external specifications can be evaluated within the moneyness range observed in earlier LO data. Online Supplement Section S4 and Table S4 report checks of the constant-volatility mapping and the number of steps in the Cox--Ross--Rubinstein tree.

This comparison tests whether the implied-volatility advantage transfers across WTI option families; it does not identify a variance risk premium.


\subsection{Independent extended Databento settlement-mark benchmark}

We assess whether the volatility-input ranking is also observed in an
independently collected AAO settlement-statistics archive. This analysis
does \emph{not} extend the earlier Curran/CL-fixing panel. Its inputs,
pricing function, and information set differ. The Databento AAO
definitions and settlement statistics cover January 2020--September
2026; admissible evaluation dates run from January 2023 through
September 2026. A fixed five-source-day sampling stride and common
contract eligibility filters give 189 valuation dates, 1,689 eligible
expiry-date groups, and 111,585 contract-date targets. Of these,
97,341 are flagged as actual settlements and 14,244 as theoretical.
All headline error comparisons below use the 97,341
actual-flagged targets.

A same-date cross-sectional put--call parity regression on a reserved
set of matched strikes supplies the risk-neutral expected terminal
average and an effective discount factor. The reserved strikes are
excluded from evaluation targets. The subsequent price is the
reduced-form Black price of a \emph{lognormal terminal arithmetic
average}, not the Curran approximation applied to a simulated path of
first-nearby CL fixes. This contemporaneous use of other AAO option
marks makes the exercise conditional on same-date market information;
it is \emph{not} an independent fully ex-ante pricing forecast.

The historical-volatility Bayesian input is estimated from lagged
changes in the parity-implied expected AAO average for the same
expiry, and thus is not a direct history of CL returns or a
risk-neutral volatility estimate. Posterior integration (PI) and
posterior-mean plug-in (PM) use the same return-proxy posterior.
Rolling and exponentially weighted proxies are evaluated on
exactly the same option targets. The unchanged lagged-implied-volatility
benchmark inverts earlier-date near-the-money AAO option marks
under the \emph{same} terminal-average pricing proxy; no target-date
implied volatility is used as a forecasting feature. All six
methods use common contract/date support.

A second exercise uses only \emph{strictly earlier} near-the-money
AAO implied volatilities to model same-expiry log-implied-volatility
innovations. Bayesian normal-inverse-gamma updating yields a
Student-\(t\) posterior predictive log-volatility distribution:
\texttt{bayes\_IV\_PI} averages option prices over 256 predictive
volatility draws; \texttt{bayes\_IV\_PM} prices at their mean.
Both are compared with the unchanged lagged-IV price on identical
options. The chronological comparison is reported separately
for 2023--2024 development dates and 2025--September 2026;
the latter dates had already been inspected during initial work and
are not an untouched confirmatory holdout for every subsequent
diagnostic.

A final \emph{distinct} calibration exercise constructs 90\% and 95\%
intervals for observed-minus-predicted \emph{settlement-mark errors},
using either strictly historical date-balanced Bayesian
normal-inverse-gamma predictive residuals or rolling empirical
residual quantiles. It uses a 63-date maximum window, requires
at least 15 earlier dates per put/call-side and maturity stratum,
and is scored by empirical coverage, interval width, and the
proper interval score. These are \emph{not} posterior credible
intervals for an Asian payoff, nor posterior price bands from
the volatility parameter in the earlier October-expiry study.
Date- and calendar-month-block resampling is used for paired
comparisons, not contract-row independence.
